
\section{Introduction}

Cellular core networks are critical infrastructure underpinning everyday communication for billions of users.
Vulnerabilities in these systems can cause large-scale service disruptions and compromise user privacy, making their security of paramount importance~\cite{son2025citesting}.
Historically, core network (CN) security has relied on physical isolation: internal interfaces between network functions operate within a trusted zone, often deployed without strong encryption or mutual authentication~\cite{3gpp33210}.
This trust model, however, is increasingly fragile. As operators migrate CNs to cloud-native deployments, the security perimeter shifts from physically isolated private infrastructure to shared cloud environments with complex network topologies~\cite{barrachina2022cloud}.
Prior work has demonstrated that attackers can reach CN interfaces through cloud misconfigurations or by abusing GTP-U tunnels to inject arbitrary traffic~\cite{zhang2025invade}.
These architectural changes expose internal interfaces to external adversaries, invalidating the implicit-trust assumption that has long underpinned core network security.

\vspace{2pt}\noindent\textbf{Vulnerabilities in CN}. Given these emerging attack surfaces, there is an urgent need for systematic approaches to uncover vulnerabilities in core network implementations. Yet, prior work has largely focused on analyzing interactions either between user equipment (UE) and CN components through the Non-Access Stratum (NAS)~\cite{garbelini2022towards, khandker2024astra, sun20255gc, dong2025corecrisis,bennett2024ransacked}, or between base stations and the CN via protocols such as the Next-Generation Application Protocol (NGAP)~\cite{sun20255gc,bennett2024ransacked} and the GPRS Tunnelling Protocol User Plane (GTP-U)~\cite{bennett2024ransacked}. To our knowledge, no systematic study has ever been done to look into security risks within the CN itself, particularly those arising from inter-component communication that may become externally exposed in cloud deployments, as discussed earlier.

On the other hand, industry has continued to report security-critical flaws affecting the internal CN~\cite{blackduck2022open5gs, nccgroup2021open5gs}. For example, Black Duck's security report~\cite{blackduck2022open5gs} reveals that Open5GS's PFCP implementation lacks proper bounds checking on IE length fields, incurring denial of service (DoS) risks. Additional vulnerabilities have surfaced in GitHub Issues, often involving open-source 5G implementations and simulators such as Open5GS, free5GC, and OpenAirInterface (OAI). Many of these findings, which are frequently uncovered using fuzzing tools such as Defensics~\cite{blackduckfuzzing}, are only isolated samples from a much larger CN security problem space, but they offer a clear glimpse into the kinds of threats these networks face. What remains missing, however, is a systematic study that connects these dots, revealing their underlying relations and, more broadly, clarifying what can go wrong in the CN.

In our research, we take a first step toward understanding CN security by analyzing vulnerabilities reported through GitHub Issues in open-source CN implementations. 
Specifically, we collected all security issues from the GitHub repositories of Open5GS and free5GC, and conducted a systematic root-cause analysis using Claude Code. We designed a Chain-of-Thought (CoT) procedure to guide automated root-cause discovery (e.g., identifying assertion-triggered crashes due to session pool exhaustion) and further clustered these flaws according to their underlying causes. Our analysis revealed a striking similarity across the reported issues. All can be grouped into three broad categories: missing validation of message syntax, missing validation of message semantics, and failure to check resource availability. Collectively, these vulnerabilities reflect a fundamental pattern of \textit{blind trust}: CN components accept internal messages as inherently trustworthy, an assumption rooted in the traditional belief that the CN is walled off from untrusted environments, but one that no longer holds in cloud-native deployments. We call them \textit{implicit trust errors}, or \textit{iTrue}. 


\vspace{2pt}\noindent\textbf{Finding iTrues}. To systematically discover iTrues, we developed a multi-agent system, \textit{iFinder}, that summarizes the root causes within each cluster, generalizes these summaries into reusable detection patterns, and applies the patterns to identify additional iTrue flaws in CN implementations. 
Achieving this required addressing a key challenge: distilling cluster-level root causes into detection patterns that are both precise and general. The patterns must capture the essential vulnerability logic while remaining broad enough to cover unseen instances across different protocols and programming languages.
To this end, we engineered prompts to emphasize language-agnostic and protocol-agnostic patterns, explicitly identifying key elements and the expected operations over them. We then use these patterns to guide the agent in scanning CN implementations and extracting iTrue candidates (Section~\ref{sec:preprocessing}). For example, \textit{iFinder} abstracts reports of missing Information Element (IE) checks into a protocol-agnostic pattern: locating downstream uses of an IE, then verifying that its presence and validity are validated earlier in message handling. We pass this pattern to the agent to drive LLM-based backward analysis and flag potential iTrues.

Moreover, naively applying these patterns to CN implementations can yield many false positives, largely due to hallucinations and incomplete reasoning when the underlying LLM lacks sufficient context. In particular, we observed that when only a local code fragment is provided, the agent may flag it as an iTrue even though the required validation is performed elsewhere (e.g., an earlier stage of the processing pipeline) outside the snippet.
At a deeper level, the problem is that protocols typically comprise multiple procedures, each spanning several states. Within a given procedure, security-critical checks may be performed in earlier states rather than in the flagged snippet, so validating an iTrue often requires analyzing the preceding states instead of a single snippet.
However, identifying the code for those preceding states is itself challenging. Feeding the entire codebase to an LLM to ``search'' for relevant checks is often both inefficient and unreliable in large, stateful CN implementations.
To mitigate this issue, we developed a novel \textit{code-specification cross-checking} technique. 
The key idea is to map an iTrue candidate to the protocol procedure it implements, as defined in the relevant specifications, and then map the entire procedure back to the codebase to recover the full execution context. This expanded, procedure-level view enables the agent to determine whether the necessary validation and resource checks are actually enforced, substantially reducing false positives when assessing whether a finding is indeed an iTrue.

The final step of \textit{iFinder} is to generate proof-of-concept (PoC) exploits to confirm the presence of an iTrue flaw. To this end, the agent follows an iterative test-and-refine loop: it deploys a simulator-based testbed to execute PoCs generated by the agent, analyzes execution logs to diagnose failures and extract actionable feedback, and then uses this feedback to revise the PoC. This process repeats until the vulnerability is successfully exploited or a testing budget is exhausted.

Running iFinder on seven prominent open-source CN implementations, we uncovered 84 previously unknown vulnerabilities in the Packet Forwarding Control Protocol (PFCP) and GPRS Tunnelling Protocol Control Plane (GTP-C) protocols. 
Among them, 83 have been confirmed by developers and 81 have been assigned CVEs.
These vulnerabilities have significant security implications, including DoS and session hijacking. 
Moreover, we evaluated these vulnerabilities on two real-world commercial 5G core networks using PoCs derived from open-source platforms, identifying two DoS vulnerabilities and one session-hijacking vulnerability. Both DoS vulnerabilities affect the same commercial 5G core network, and we have received a CVE assignment for one of them (CVE-2026-8232).
The session-hijacking flaw was found in both commercial 5G core networks: one vendor has since fixed the issue, with a CVE assigned (CVE-2026-8233), while the other is still in the remediation process. At a high level, this flaw allows an attacker to inject a PFCP Session Modification Request, causing the User Plane Function (UPF) to forward the victim UE’s uplink traffic to the attacker. We include this vulnerability as a case study in Section~\ref{sec: case studies}.

\noindent \textbf{Contributions.} We list our following contributions below:

\vspace{2pt}\noindent$\bullet$ \textit{New findings}. We present the first systematic root-cause analysis of vulnerabilities reported in open-source CN implementations, showing that they largely arise from \textit{implicit trust} among internal CN components.

\vspace{2pt}\noindent$\bullet$ \textit{New intelligent techniques}. We developed \textit{iFinder}, the first multi-agent system for automated vulnerability discovery in carrier core networks. Our techniques address key technical barriers, and an ablation study demonstrates their effectiveness in reducing hallucinations.

\vspace{2pt}\noindent$\bullet$ \textit{Implementation and evaluation}. We implemented \textit{iFinder} and evaluated it on seven widely used open-source CN implementations, uncovering \textbf{84} previously unknown vulnerabilities and \textbf{81} assigned CVEs, including a session-hijacking flaw confirmed on real-world commercial 5G core networks.

